\documentclass{article}
\usepackage[UTF8]{ctex}
\usepackage{amsmath,amssymb,amsfonts}
\usepackage{enumitem}
\usepackage{geometry}
\geometry{a4paper, margin=1in}
% === 错题本排版优化（不参与题目标识） ===
\setlength{\parindent}{0pt}
\setlength{\parskip}{0.35em plus 0.10em minus 0.08em}
\setlength{\emergencystretch}{3em}
\setlist{itemsep=0.25em, topsep=0.40em, parsep=0pt}
\allowdisplaybreaks[2]
% === 错题本排版优化结束 ===
\begin{document}
\section*{南京师范大学2024年数学分析真题}
\section*{91．南京师范大学2024年数学分析试卷}

\subsection{题目1}
一．计算题．每题5分，共20分．

\begin{itemize}
  \item[（1）] 求极限 \(\lim _{n \rightarrow \infty} \tan ^{n}\left(\frac{\pi}{4}+\frac{1}{n}\right)\) ．
  \item[（2）] 已知 \(\lim _{x \rightarrow-\infty}\left(f(x)-f^{\prime}(x)\right)=2023\) ，求 \(\lim _{x \rightarrow-\infty} f(x)\) ．
  \item[（3）] 求定积分 \(\int_{0}^{a} x^{2} \sqrt{\frac{a-x}{a+x}} \mathrm{~d} x(a>0)\) ．
  \item[（4）] 交换积分顺序 \(\int_{\frac{1}{2}}^{\frac{1}{\sqrt{2}}} \mathrm{~d} x \int_{\frac{1}{2}}^{x} f(x, y) \mathrm{d} y+\int_{\frac{1}{\sqrt{2}}}^{1} \mathrm{~d} x \int_{x^{2}}^{x} f(x, y) \mathrm{d} y\) 。
\end{itemize}

\subsection{题目2}
二．（15 分）设级数 \(\sum_{n=1}^{\infty} a_{n}\) 收敛于 \(a\) 。

\begin{itemize}
  \item[（1）] 求极限 \(\lim _{n \rightarrow \infty} \frac{n a_{1}+(n-1) a_{2}+\cdots+a_{n}}{n}\) ．
  \item[（2）] 证明： \(\lim _{n \rightarrow \infty} \frac{a_{1}+2 a_{2}+\cdots+n a_{n}}{n}=0\) ．
\end{itemize}

\subsection{题目3}
三．（15 分）设 \(f(x)\) 在 \(x=0\) 处连续，且对一切 \(x, y \in \mathbb{R}\) ，有 \(f(x+y)=f(x)+f(y)\) 。证明：\(f(x)\) 在\(\mathbb{R}\) 上都续且 \(f(x)=f(1) x\) ．

\subsection{题目4}
四．（15 分）设 \(f(x)\) 在 \([a, b]\) 上二次可微，\(f^{\prime}\left(\frac{a+b}{2}\right)=0\) ，证明：存在 \(\xi \in(a, b)\) ，使得

\[
\left|f^{\prime \prime}(\xi)\right| \geq \frac{4}{(b-a)^{2}}|f(b)-f(a)|
\]

\subsection{题目5}
五．（15 分）解答如下问题：

\begin{itemize}
  \item[（1）] 证明：若 \(f(x), g(x)\) 均在 \([a, b]\) 上可积，则有
\[
\left(\int_{a}^{b} f(x) g(x) \mathrm{d} x\right)^{2} \leq \int_{a}^{b} f^{2}(x) \mathrm{d} x \int_{a}^{b} g^{2}(x) \mathrm{d} x
\]
  \item[（2）] 证明：若 \(f^{\prime}(x)\) 在 \([0,1]\) 上连续，\(f(0)=0\) ，则有
\[
\int_{0}^{1} f^{2}(x) \mathrm{d} x \leq \frac{1}{2} \int_{0}^{1}\left[f^{\prime}(x)\right]^{2} \mathrm{~d} x .
\]
\end{itemize}

\subsection{题目6}
六。（10分）设函数 \(f(x)\) 在 \([a,+\infty)\) 上单调递减，且 \(f(x)>0, x \in[a,+\infty)\) ，证明：反常积分\(\int_{a}^{+\infty} f(x) \mathrm{d} x\) 与 \(\int_{a}^{+\infty} f(x) \sin ^{2} x \mathrm{~d} x\) 敛散性相同．

\subsection{题目7}
七．（20 分）解答如下问题：\\
（1）已知 \(f(x, y)\) 在点 \((0,0)\) 的某邻域内连续，且

\[
\lim _{(x, y) \rightarrow(0,0)} \frac{f(x, y)-x y^{2}}{1-\cos \sqrt{x^{2}+y^{2}}}=1 .
\]

试问：（ 0,0 ）是否为 \(f(x, y)\) 的驻点（或稳定点）？进一步讨论（ 0,0 ）是否为 \(f(x, y)\) 的极值点？若是是极大值点还是极小值点？若不是，请说明理由。\\

\subsection{题目7}
原题第（2）问：设 \(u=u(x, y, z), x=r \sin \varphi \cos \theta, y=r \sin \varphi \sin \theta, z=r \cos \varphi\) ，若 \(x u_{x}+y u_{y}+z u_{z}=0\) ，证明 \(u\) （见为 \(\theta\) 与 \(\varphi\) 的函数。

\subsection{题目8}
八．（15 分）讨论级数 \(\sum_{n=1}^{\infty} \frac{(-\alpha)^{n} n!}{n^{n}}(\alpha>0)\) 的敛散性，若收敛，请指出是条件收敛还是绝对收敛。

\subsection{题目9}
九．（15 分）计算曲面积分

\[
\iint_{S}\left(x^{2} \cos \alpha+y^{2} \cos \beta+z^{2} \cos \gamma\right) \mathrm{d} S
\]

其中 \(S\) 为圆锥曲面 \(x^{2}+y^{2}=z^{2}\) 上位于 \(0 \leq z \leq h\) 的一部分，其法向量恒与 \(z\) 轴正向相交成锐角，\(\cos \alpha, \cos \beta, \cos \gamma\) 为法向量的方向余弦。

\subsection{题目10}
十．（10 分）证明： \(\int_{0}^{+\infty} \frac{\cos x^{2}}{x^{y}} \mathrm{~d} x\) 关于参变量 \(y\) 在 \((-1,1)\) 内内闭一致收敛．
\end{document}
